\documentclass[11pt]{article}

% ── Packages ──────────────────────────────────────────────────────────────────
\usepackage{neurips_2024}
\usepackage[utf8]{inputenc}
\usepackage[T1]{fontenc}
\usepackage{amsmath}
\usepackage{amssymb}
\usepackage{amsfonts}
\usepackage{amsthm}
\usepackage{algorithm}
\usepackage{algorithmic}
\usepackage{booktabs}
\usepackage{enumitem}
\usepackage{listings}
\usepackage{xcolor}
\usepackage{pifont}
\usepackage{hyperref}
\usepackage{url}
\usepackage{natbib}

% ── Theorem environments ─────────────────────────────────────────────────────
\newtheorem{theorem}{Theorem}
\newtheorem{lemma}[theorem]{Lemma}
\newtheorem{proposition}[theorem]{Proposition}
\newtheorem{corollary}[theorem]{Corollary}
\newtheorem{definition}{Definition}
\newtheorem{remark}{Remark}

% ── Listings configuration ───────────────────────────────────────────────────
\lstset{
  language=Python,
  basicstyle=\ttfamily\small,
  keywordstyle=\color{blue}\bfseries,
  commentstyle=\color{gray}\itshape,
  stringstyle=\color{red!70!black},
  numberstyle=\tiny\color{gray},
  numbers=left,
  numbersep=5pt,
  breaklines=true,
  showstringspaces=false,
  frame=single,
  rulecolor=\color{black!30},
  backgroundcolor=\color{gray!5},
  tabsize=4,
  captionpos=b
}

% ── Convenience macros ───────────────────────────────────────────────────────
\newcommand{\cmark}{\ding{51}}
\newcommand{\xmark}{\ding{55}}

% ══════════════════════════════════════════════════════════════════════════════
\title{TSAM: Typed Serving with Tenant-Scoped Attention Masking\\
for Provably Secure Multi-Tenant LLM Inference}

\author{Anonymous Authors\\NeurIPS 2024 Submission}

\begin{document}
\maketitle

% ══════════════════════════════════════════════════════════════════════════════
% ABSTRACT
% ══════════════════════════════════════════════════════════════════════════════
\begin{abstract}
Multi-tenant serving of large language models (LLMs) is economically
essential---GPU costs make dedicated instances per tenant prohibitive at
scale---yet sharing a single model instance across tenants introduces
cross-tenant information flow through the shared key--value (KV) cache.
Despite significant recent interest, \emph{no existing approach
simultaneously achieves} (1)~provably zero information leakage between
tenants, (2)~unlimited tenant scalability, and (3)~compatibility with
prefix sharing, the principal memory optimization in production serving.
Orthogonal subspace projection~\citep{li2024orthogonal} is limited to at
most $d_{\mathrm{model}}/d_{\mathrm{sub}} \leq 16$ tenants and
fundamentally destroys prefix sharing because each tenant's keys occupy a
private subspace.  Process-level isolation preserves correctness but
incurs 5--10$\times$ cost overhead, eliminating the economic rationale for
shared serving.

We present \textbf{TSAM} (\textbf{T}yped \textbf{S}erving with
Tenant-Scoped \textbf{A}ttention \textbf{M}asking), a lightweight
modification to PagedAttention~\citep{kwon2023vllm} that closes this gap.
TSAM assigns each KV-cache page a \emph{type tag} drawn from a small set
$\{\texttt{shared}, \texttt{private}(t)\}$ and enforces a single
hard-masking conditional in the attention kernel: a query belonging to
tenant~$t_i$ may attend to a page only if that page is tagged
\texttt{shared} or \texttt{private}($t_i$).  We prove that the resulting
mutual information satisfies
$I\!\bigl(\mathrm{output}(q_i);\;\mathrm{KV}_{\mathrm{private}}(t_j)\bigr)
= 0$ for all $i \neq j$---a \emph{structural zero}, not a statistical
approximation.  Empirically, TSAM achieves \textbf{0.0 bits} of measured
cross-tenant leakage, adds \textbf{$<$1\%} latency overhead to
FlashAttention-2~\citep{dao2022flashattention} kernels, scales to an
\textbf{unlimited} number of tenants, and delivers $\geq$\textbf{95\%}
reduction in timing side-channel distinguishability via an optional
constant-time padding defense.

To the best of our knowledge, TSAM is the \emph{first} system to
simultaneously achieve all three properties: provable zero leakage,
arbitrary tenant scale, and full prefix sharing compatibility.
\end{abstract}

% ══════════════════════════════════════════════════════════════════════════════
% SECTION 1 — INTRODUCTION
% ══════════════════════════════════════════════════════════════════════════════
\section{Introduction}
\label{sec:introduction}

Large language models (LLMs) such as GPT-4~\citep{brown2020language,
achiam2023gpt4} power an expanding array of commercial applications---from
customer-facing chatbots to enterprise document analysis---that serve
thousands of tenants from a shared infrastructure.  The economic case is
compelling: a single A100 or H100 GPU can cost \$2--4/hour, and dedicating
one instance per tenant is infeasible beyond a handful of customers.
Systems such as vLLM~\citep{kwon2023vllm} address this through
\emph{PagedAttention}, which manages the KV-cache in fixed-size pages and
enables \emph{continuous batching} of requests from different tenants
within the same forward pass.  This design yields order-of-magnitude
improvements in GPU utilization and throughput.

However, batching requests from distinct tenants into a single attention
computation creates a subtle but critical vulnerability: tokens from
tenant~$t_i$ can, in principle, attend to KV-cache entries produced by
tenant~$t_j$'s private context.  Recent work has shown that LLMs can
memorize and regurgitate training data~\citep{carlini2021extracting}, and
analogous risks arise at inference time when private KV-cache pages are
accessible to cross-tenant queries.  The attention
mechanism~\citep{vaswani2017attention} computes
$\mathrm{softmax}(QK^\top/\sqrt{d_k})\,V$ over \emph{all} keys in the
batch; without explicit isolation, every query has a nonzero attention
weight on every key.

\paragraph{The gap.}
Despite the critical importance of multi-tenant LLM security, no existing
approach simultaneously achieves (1)~provably zero information leakage
between tenants, (2)~arbitrary tenant scalability, and (3)~compatibility
with prefix sharing---the technique that avoids redundant computation of
common system prompts.  This gap is a fundamental barrier to deploying
shared LLM serving in regulated industries (healthcare, finance, legal)
where formal isolation guarantees are a prerequisite.

\paragraph{Limitations of existing approaches.}
\emph{Process isolation} runs each tenant in a separate serving instance.
While trivially correct, this approach incurs 5--10$\times$ cost overhead
due to duplicated model weights and KV-cache, entirely negating the
economic advantage of shared serving.
\emph{Orthogonal subspace projection}~\citep{li2024orthogonal} maps each
tenant's keys into a mutually orthogonal subspace of dimension
$d_{\mathrm{sub}}$, ensuring that cross-tenant dot products vanish.
However, the number of orthogonal subspaces is bounded by
$d_{\mathrm{model}} / d_{\mathrm{sub}}$; for typical models with
$d_{\mathrm{model}} = 4096$ and $d_{\mathrm{sub}} = 256$, this yields a
hard ceiling of \textbf{16 tenants}.  Moreover, because each tenant's keys
are rotated into a private subspace, prefix sharing is fundamentally
destroyed: the same prompt produces different key vectors for different
tenants, and no two tenants can share a cache page.
\emph{Ad hoc attention masks} have been explored in concurrent
work~\citep{gao2024attentionguard}, but these approaches lack formal
information-flow proofs and do not address timing side channels.

\paragraph{Our approach.}
We present \textbf{TSAM} (Typed Serving with Tenant-Scoped Attention
Masking), a principled and lightweight extension to PagedAttention that
resolves all three limitations simultaneously.  TSAM introduces a
single-bit \emph{type tag} per KV-cache page---\texttt{shared} or
\texttt{private}($t$)---and enforces a hard-masking rule in the attention
kernel: query $q$ from tenant $t_i$ may attend to page $p$ if and only if
$\mathrm{type}(p) = \texttt{shared}$ or $\mathrm{owner}(p) = t_i$.  This
conditional adds fewer than 10 instructions per page in the inner loop of
FlashAttention-2 and yields a \emph{structural zero} in mutual
information:
\begin{equation}
  I\!\bigl(\mathrm{output}(q_i);\;\mathrm{KV}_{\mathrm{private}}(t_j)\bigr) = 0
  \quad \forall\; i \neq j.
  \label{eq:structural_zero}
\end{equation}
Because the mask is enforced at the page level rather than the subspace
level, tenant count is bounded only by available memory, and shared-prefix
pages remain tagged \texttt{shared}---preserving prefix sharing without
modification.

\paragraph{Contributions.}
We make the following contributions:
\begin{enumerate}[leftmargin=*,itemsep=2pt]
\item \textbf{Typed KV-cache pages.} We introduce a page-level type
  system for the KV-cache that cleanly separates shared and
  tenant-private state within the existing PagedAttention framework
  (\S\ref{sec:method}).

\item \textbf{Tenant-scoped attention masking.} We design a hard-masking
  rule that is enforced inside the attention kernel, achieving provably
  zero cross-tenant information flow with a single conditional per page
  (\S\ref{sec:method}).

\item \textbf{Formal security proof.} We prove that TSAM satisfies
  noninterference~\citep{goguen1982security} under the lattice model of
  information flow~\citep{denning1976lattice}, yielding $I = 0$ as a
  structural invariant rather than a statistical bound
  (\S\ref{sec:formal}).

\item \textbf{Timing side-channel defense.} We present a constant-time
  padding mechanism that eliminates sequence-length-dependent timing
  variation, reducing timing distinguishability by $\geq$95\%
  (\S\ref{sec:timing}).

\item \textbf{Comprehensive evaluation.} We evaluate TSAM on
  Llama-2-7B and Llama-2-13B under realistic multi-tenant workloads,
  demonstrating 0.0 bits leakage, $<$1\% latency overhead, and linear
  scaling to hundreds of concurrent tenants (\S\ref{sec:experiments}).
\end{enumerate}

\paragraph{Comparison with existing approaches.}
Table~\ref{tab:comparison} summarizes the landscape.  TSAM is the only
approach that achieves all four desiderata.

\begin{table}[t]
\centering
\caption{Comparison of multi-tenant LLM isolation approaches.}
\label{tab:comparison}
\vspace{4pt}
\small
\begin{tabular}{@{}lcccc@{}}
\toprule
\textbf{Approach} & \textbf{Max Tenants} & \textbf{Prefix Sharing} & \textbf{Formal Guarantee} & \textbf{Overhead} \\
\midrule
Process Isolation          & Unlimited & \xmark & Trivial       & 5--10$\times$ \\
Orth.\ Projection~\citep{li2024orthogonal} & $\leq 16$ & \xmark & $\epsilon$-approx & Moderate \\
Ad hoc Masks~\citep{gao2024attentionguard}  & Unlimited & \cmark & None          & Low \\
\textbf{TSAM (ours)}      & \textbf{Unlimited} & \cmark & \textbf{$I{=}0$ (structural)} & \textbf{$<$1\%} \\
\bottomrule
\end{tabular}
\end{table}

% ══════════════════════════════════════════════════════════════════════════════
% SECTION 2 — BACKGROUND AND RELATED WORK
% ══════════════════════════════════════════════════════════════════════════════
\section{Background and Related Work}
\label{sec:background}

\subsection{Multi-Tenant LLM Serving}
\label{sec:bg_serving}

Modern LLM serving systems process requests from many tenants on shared
GPU infrastructure.  At the core of efficient serving lies the
\emph{key--value (KV) cache}: during autoregressive generation, the key
and value projections for all previously generated tokens are stored so
that attention need not be recomputed from scratch at each decoding step.
For a model with $L$ layers, $H$ attention heads, and head dimension
$d_h$, the KV-cache for a single sequence of length $n$ requires
$2 L H d_h n$ elements, consuming tens of gigabytes for long contexts.

\textbf{PagedAttention}~\citep{kwon2023vllm} addresses KV-cache memory
fragmentation by managing the cache in fixed-size \emph{pages} (typically
16 tokens each), analogous to virtual memory in operating systems.  Pages
are allocated on demand and can be shared across sequences via
copy-on-write, enabling \emph{prefix sharing}: when multiple requests
begin with the same system prompt, their common prefix pages are stored
once and referenced by all.

\textbf{Continuous batching}~\citep{yu2022orca} further improves
throughput by dynamically inserting and removing requests from the GPU
batch at each iteration, rather than waiting for an entire batch to
complete.  Together, PagedAttention and continuous batching allow a single
GPU to serve hundreds of concurrent sequences with high utilization.

\textbf{Prefix sharing} is especially important in multi-tenant settings
where many tenants share a common system prompt (e.g., ``You are a helpful
assistant'').  Without prefix sharing, the KV-cache for shared prompts
must be duplicated per tenant, wasting memory proportional to the number
of tenants times the prompt length.  Any isolation mechanism that destroys
prefix sharing therefore imposes a significant and scaling cost penalty.

\subsection{Information Flow Security}
\label{sec:bg_infoflow}

The formal study of information flow in computing systems originates with
Denning's \emph{lattice model}~\citep{denning1976lattice}, which assigns
security labels to data objects and defines permitted information flows via
a partial order on labels.  A system satisfies \emph{secure information
flow} if, for every pair of labels $\ell_1 \not\leq \ell_2$, no
information flows from objects at level $\ell_1$ to observers at level
$\ell_2$.

Goguen and Meseguer~\citep{goguen1982security} formalized this intuition
as \emph{noninterference}: a system is noninterfering with respect to a
high-security domain $H$ and a low-security domain $L$ if the outputs
visible to $L$ are identical regardless of the inputs from $H$.  Formally,
for a system function $f$ with inputs partitioned into $H$ and $L$
components:
\begin{equation}
  f(h_1, l) = f(h_2, l) \quad \forall\; h_1, h_2 \in H,\; \forall\; l \in L.
  \label{eq:noninterference}
\end{equation}
In the multi-tenant LLM setting, we treat each tenant's private data as a
separate high-security domain and require that the outputs for tenant
$t_i$ are independent of the private inputs of any other tenant $t_j$
($j \neq i$).

An equivalent information-theoretic characterization uses mutual
information~\citep{shannon1948mathematical, cover2006elements}: a system
achieves perfect isolation if
$I(\mathrm{output}_i;\;\mathrm{input}_j) = 0$ for all $i \neq j$.  When
this quantity is exactly zero (not merely bounded by $\epsilon$), we call
it a \emph{structural zero}---it holds by construction rather than by
parameter tuning or statistical estimation.

\subsection{Cache Side-Channel Attacks}
\label{sec:bg_sidechannel}

Hardware cache side-channel attacks exploit timing differences in memory
access patterns to infer secret data.
Osvik~et~al.~\citep{osvik2006cache} demonstrated \emph{access-driven}
attacks on AES, where the attacker observes which cache lines are evicted
during encryption.  Yarom and Falkner~\citep{yarom2014flush} introduced
\textsc{Flush+Reload}, a high-resolution attack exploiting shared memory
pages in last-level caches to achieve cross-core information leakage.

In the context of LLM serving, analogous side channels arise through the
\emph{KV-cache}.  When multiple tenants share a PagedAttention cache pool,
the set of pages allocated to a tenant, the timing of page allocations and
evictions, and the sequence lengths (which determine the number of pages
accessed during attention) can all leak information.  Specifically:
\begin{itemize}[leftmargin=*,itemsep=2pt]
  \item \textbf{Attention weight leakage.} Without masking, softmax
    attention assigns nonzero weight to every key in the batch, including
    keys from other tenants' private contexts.  Even small attention
    weights create a nonzero information channel.
  \item \textbf{Timing side channels.} The latency of attention
    computation depends on the total number of KV-cache pages accessed.
    An adversarial tenant can infer the sequence length of co-located
    tenants by measuring response latency, potentially revealing
    confidential information about the nature of their queries.
  \item \textbf{Cache collision attacks.} In shared GPU memory, two
    tenants accessing the same physical cache lines can observe
    interference patterns, analogous to classical cache collision attacks
    in CPUs.
\end{itemize}
TSAM addresses the first two channels directly: hard masking eliminates
attention weight leakage, and optional constant-time padding neutralizes
timing variation.  GPU-level cache collision attacks are orthogonal and
require hardware-level mitigations beyond the scope of this work.

\subsection{Existing Isolation Mechanisms}
\label{sec:bg_isolation}

We survey three categories of existing approaches to multi-tenant LLM
isolation.

\paragraph{Process isolation.}
The most straightforward approach runs a separate model instance per
tenant.  Each tenant's KV-cache is physically isolated, trivially
guaranteeing noninterference.  However, this approach requires duplicating
the full model weights ($\sim$14\,GB for a 7B-parameter model in FP16)
and KV-cache per tenant, yielding 5--10$\times$ memory overhead and
proportionally reduced throughput.  For $N$ tenants, the total memory
scales as $O(N \cdot (W + C))$ where $W$ is the model weight size and $C$
is the per-tenant cache size, compared to $O(W + N \cdot C)$ for shared
serving.  Process isolation is therefore economically viable only for a
small number of high-value tenants.

\paragraph{Orthogonal subspace projection.}
Li~et~al.~\citep{li2024orthogonal} propose projecting each tenant's key
vectors into mutually orthogonal subspaces of the attention head
dimension.  Given head dimension $d_h$, each tenant is assigned a
subspace of dimension $d_{\mathrm{sub}}$, and the key projection is
followed by a projection matrix $P_t \in \mathbb{R}^{d_h \times d_h}$
that zeros out all components outside tenant $t$'s subspace.  Since the
subspaces are orthogonal, cross-tenant dot products vanish:
\begin{equation}
  q_i^\top P_{t_i}^\top P_{t_j} k_j = q_i^\top (P_{t_i} P_{t_j})^\top k_j = 0
  \quad \text{for } i \neq j,
  \label{eq:orthogonal}
\end{equation}
because $P_{t_i} P_{t_j} = 0$ for orthogonal projections onto disjoint
subspaces.  This approach has two critical limitations:
\begin{enumerate}[leftmargin=*,itemsep=2pt]
  \item \textbf{Tenant scalability.} The maximum number of tenants is
    $d_h / d_{\mathrm{sub}}$.  For Llama-2's head dimension $d_h = 128$
    and a minimum practical subspace dimension $d_{\mathrm{sub}} = 8$,
    this yields at most \textbf{16 tenants}.  Reducing $d_{\mathrm{sub}}$
    further degrades model quality because the effective attention
    dimension per tenant shrinks.
  \item \textbf{Prefix sharing destruction.} Because each tenant's keys
    are multiplied by a different projection matrix $P_t$, the same input
    token produces different key vectors for different tenants.
    Consequently, no two tenants can share KV-cache pages for a common
    prefix---prefix sharing is \emph{fundamentally incompatible} with
    subspace projection.
\end{enumerate}

\paragraph{Ad hoc attention masks.}
Several concurrent works~\citep{gao2024attentionguard, rajput2024inference}
propose applying attention masks to prevent cross-tenant attention.  While
directionally similar to TSAM, these approaches differ in important ways:
they operate at the token level rather than the page level (incurring
higher metadata overhead), they do not provide formal information-flow
proofs or noninterference guarantees, and they do not address timing side
channels.  TSAM's contribution is to show that page-level typing, combined
with a formal security model grounded in the lattice framework of
Denning~\citep{denning1976lattice}, yields a provably correct and
practically efficient isolation mechanism.